\begin{textsample}{POS Dim 4 – 2009\_c6 – Score -3.00 – 2009\_c6/enade\_2009\_com\_soc\_cine\_3a63805496.txt}  \label{ex:f4_pos_322}
Capítulo 3 Análise Técnica da Prova Este capítulo tem por objetivo apresentar o desempenho dos estudantes de COMUNICAÇÃO SOCIAL – CINEMA no ENADE/2009.
Para isso, foram calculadas as estatísticas básicas da prova como um todo, bem como as estatísticas das questões relacionadas à Formação Geral e ao Componente Específico.
Nas Tabelas, são apresentadas as seguintes estatísticas : total da população de presentes, média, erro-padrão da média, desvio-padrão, mediana, nota máxima, nota mínima e coeficiente de assimetria ( ver definição em 1.4.3 ).
Nas estatísticas das questões discursivas, foram detalhadas também as notas nulas : em branco, por protesto ou por insuficiência.
As estatísticas apresentadas neste capítulo contemplam, separadamente, os ingressantes, os concluintes e o total de estudantes, e foram calculadas tendo -se em vista as seguintes agregações : ( a ) Região Geográfica e País e ( b ) a Categoria Administrativa e a Organização Acadêmica.
As notas brutas das questões discursivas em cada uma das partes, Formação Geral/Componente Específico, foram calculadas a partir da média simples das notas de cada uma das questões correspondentes.
A primeira questão discursiva de Formação Geral foi desconsiderada e, portanto, a nota bruta referente às questões discursivas de Formação Geral foi igual à nota da questão mantida, a segunda.
No Componente Específico, a nota bruta correspondeu à média das notas nas três questões.
As notas brutas das questões objetivas de cada uma das partes, Formação Geral/Componente Específico, foram calculadas a partir dos acertos dos estudantes considerando -se uma regra de três para a conversão para a escala de 0 a 100, descontando -se aquelas questões que foram desconsideradas pela comissão ou que não passaram pelo teste do índice de discriminação ponto-bisserial ( ver seção 1.4.6 para a descrição do índice ).
As questões com índices de discriminação inferior a 0,2 foram excluídas do cômputo da nota do estudante.
As notas brutas de Formação Geral e do Componente Específico para cada um dos estudantes foram calculadas como médias ponderadas das suas notas nas questões objetivas e discursivas, ponderação diferente para cada parte : C\_FG = C\_FG\textasciicircum{}OBJ × 0,6 + C\_FG\textasciicircum{}DIS × 0,4 C\_CE = C\_CE\textasciicircum{}OBJ × 0,85 + C\_CE\textasciicircum{}DIS × 0,15 onde C\_FG é a nota bruta da Formação Geral, C\_FG\textasciicircum{}OBJ é a nota bruta nas questões objetivas da Formação Geral, C\_FG\textasciicircum{}DIS é a nota bruta nas questões objetivas da Formação Geral, C\_CE é a nota bruta do Componente Específico, C\_CE\textasciicircum{}OBJ é a nota bruta nas questões objetivas do Componente Específico, e C\_CE\textasciicircum{}DIS é a nota bruta nas questões objetivas do Componente Específico.
A nota bruta da prova, C\_GER, por sua vez, é também uma média ponderada da nota de cada uma das partes, Formação Geral/Componente Específico : C\_GER = 0,25 × C\_FG + 0,75 × C\_CE. 3.1 Estatísticas Básicas da Prova 3.1.1 Estatísticas Básicas Gerais A Tabela 3.1 apresenta as estatísticas básicas da prova por grupo de estudantes.
O total de estudantes inscritos para a prova foi de 967.
Destes, 74,7 \% compareceram, sendo que a presença foi maior entre concluintes ( 83,2 \% ) do que entre ingressantes ( 67,0 \% ).
A média geral da prova foi de 44,53, sendo a média dos ingressantes menor ( 44,51 ) em relação aos concluintes ( 44,54 ).
O desvio-padrão geral foi de 22,4, sendo o do grupo de ingressantes menor ( 19,3 ), do que o do grupo dos concluintes ( 25,0 ), aqueles possuíram uma distribuição mais homogênea nas notas.
A nota máxima foi 88,8, obtida por pelo menos um concluinte ao passo que a maior nota obtida por um ingressante foi 78,2.
Os coeficientes de assimetria para os dois grupos considerados, o de concluintes e o de ingressantes, bem como para a população como um todo foram todos negativos, indicando uma maior concentração à direita da distribuição e um maior espalhamento à esquerda.
Uma comparação entre o desempenho de ingressantes e concluintes pode ser realizada a partir dos dados apresentados no Gráfico 3.1.
Cumpre destacar que as notas dos concluintes, de um modo geral, foram mais elevadas do que as dos ingressantes, cuja distribuição de notas está à esquerda dos primeiros.
A ideia de melhor desempenho é reforçada pelas notas medianas dos dois grupos : 52,6 entre os concluintes e 48,9 entre os ingressantes.
Nas faixas de notas até 60 pontos, o percentual de ingressantes é maior que o dos concluintes, situação esta que se inverte para as faixas de notas superiores a 60 pontos.
A exceção é o primeiro intervalo de 0 a 10 que inclui uma quantidade maior de concluintes do que de ingressantes.
Ainda assim, aproximadamente 46,3 \% dos ingressantes situaram -se nas faixas de notas superiores a 50,0 pontos, por oposição a 54,1 \% dos concluintes.
O Gráfico 3.2 apresenta informações referentes ao desempenho geral de ingressantes e concluintes, comparando os resultados considerando a Região, Categoria Administrativa e Organização Acadêmica.
Levando -se em conta as notas médias dos estudantes em cada Região, observa -se, nos dois grupos, há uma variação expressiva entre as Regiões.
As notas médias mais elevadas entre os concluintes ocorreram nas Regiões Centro-Oeste ( 58,17 ) e Sudeste ( 45,41 ), ambas acima da média nacional ( 44,54 ).
Em relação aos ingressantes, as médias mais altas foram obtidas nas Regiões Nordeste ( 54,60 ), Sudeste ( 47,77 ) e Centro-Oeste ( 45,31 ) ; nas demais Regiões a média foi inferior à média nacional ( 44,51 ).
Na Região Sul foi obtida a média de notas mais baixa entre os concluintes ( 38,15 ) e os ingressantes ( 32,13 ) O Sudeste apresentou a média dos ingressantes superior à dos concluintes, fato devido ao maior número de provas com nota zero entre os concluintes.
Cursos nesta área só foram oferecidos em instituições federais, estaduais e particulares.
Levando -se em conta os agrupamentos dos estudantes em categorias administrativas, observa -se que a pontuação média mais elevada entre os concluintes e ingressantes foi encontrada em instituições de origem particular ( 53,65 ) – ambas acima da média nacional, se observa que a pontuação média mais elevada entre os ingressantes foi encontrada em instituições de origem federal ( 56,46 ) – também acima da média nacional.
Tendo como foco as notas médias dos estudantes agrupados por Organização Acadêmica, notamos que Centros Universitários apresentaram as maiores médias entre os concluintes ( 51,78 ) e Faculdades apresentaram as maiores médias entre os ingressantes ( 50,65 ).
Ao passo que as piores notas médias foram observadas nas Faculdades para concluintes ( 47,33 ) e nos Centros Universitários para os ingressantes ( 35,40 ). 3.1.2 Estatísticas Básicas de Formação Geral A Tabela 3.2 apresenta as estatísticas básicas de Formação Geral por grupo de estudantes.
O número de estudantes inscritos no exame foi de 967.
Destes, 74,7 \% compareceram, sendo que a presença foi maior entre concluintes ( 83,2 \% ) do que entre ingressantes ( 67,0 \% ).
A média geral da parte da prova de Formação Geral foi de 41,13, sendo que os ingressantes obtiveram média mais alta, 42,98, que os concluintes, cuja média foi 39,48.
O desvio-padrão geral foi de 24,4, sendo o do grupo de ingressantes menor ( 23,2 ), do que o do grupo de concluintes ( 25,4 ), aqueles possuíram uma distribuição mais homogênea nas notas.
A nota máxima foi 91,4, obtida por estudantes de ambos os grupos.
Para esta parte assim como para a nota geral, os coeficientes de assimetria são negativos para todos os grupos.
Como, porém, os valores são bem perto de zero estes indicaram uma ligeira inclinação para a direita.
O Gráfico 3.3 permite a comparação entre concluintes e ingressantes relativamente ao desempenho na parte da prova que aborda a Formação Geral.
Nesta parte, se excluídas as provas com notas zero, tal como na prova como num todo, também os concluintes apresentaram desempenho ligeiramente inferior, embora seja perceptível um equilíbrio.
A mediana dos concluintes foi igual à dos ingressantes, sendo a média dos concluintes inferior à dos ingressantes.
Dos ingressantes, aproximadamente 40 \% alcançaram nota superior a 50 ; dentre os concluintes, o percentual dos estudantes nesta faixa de notas foi muito semelhante, 38,1 \%.
No Gráfico 3.4 são apresentadas as informações referentes ao desempenho dos ingressantes e dos concluintes na parte de Formação Geral, em diferentes agregações como Região, Categoria Administrativa e Organização Acadêmica.
Considerando -se as notas médias dos estudantes segundo cada Região, observou -se que, tanto no grupo dos ingressantes quanto no grupo dos concluintes, houve diferenças marcantes.
As médias variam, entre os concluintes, de 35,05 na Região Sul a 47,69 na Região Centro-Oeste.
Entre os ingressantes, os valores foram de 31,68 na Região Sul a 55,56 na Região Nordeste.
A maior dispersão de notas dos estudantes de uma dada Região foi encontrada no Sul para os concluintes e no Nordeste para os ingressantes.
Nas Regiões Sul e Centro-Oeste, os concluintes apresentaram médias maiores do que os ingressantes, a situação inversa aconteceu no Sudeste.
A Região Nordeste não apresentou concluintes e a Região Norte nem concluintes nem ingressantes.
Tendo como foco as categorias administrativas, destacou -se que somente as instituições federais, estaduais e particulares estavam representadas, sendo que as notas médias das federais foram superiores para ingressantes ( 57,49 ).
Por outro lado, as notas médias das particulares foram superiores para concluintes ( 48,75 ).
Os estudantes das instituições particulares se apresentam mais homogêneos, já que o desvio padrão para concluintes e ingressantes é o menor dos valores encontrados.
Considerando -se o tipo de Organização Acadêmica, os Centros Universitários apresentaram o maior valor para concluintes ( 44,80 ), não apresentando ingressantes.
E as Universidades apresentaram o maior valor para ingressantes ( 46,84 ), sendo este inclusive superior a nota dos concluintes.
Os estudantes das Faculdades se apresentam menos homogêneos, já que o desvio padrão é o maior encontrado em ambos os grupos.
Os Centros Federais de Educação Tecnológica e os Institutos Federais de Educação, Ciência e Tecnologia não estavam representados nesta área. 3.1.3 Estatísticas Básicas do Componente Específico A Tabela 3.3 apresenta as estatísticas básicas referentes ao Componente Específico da prova, tendo sua média geral, de 45,70, inferior ao desempenho da prova de COMUNICAÇÃO SOCIAL – CINEMA como um todo, cuja média foi 44,53, como apresentado na Tabela 3.1, que por sua vez foi inferior ao desempenho na Formação Geral, que teve média igual a 41,13, conforme mostrado na Tabela 3.2.
Os concluintes apresentaram um desempenho médio ( 46,27 ) superior ao dos ingressantes ( 45,07 ), havendo ainda indícios de maior homogeneidade entre as notas médias dos ingressantes do que entre as dos concluintes, como sugeriu o desvio padrão, menor no grupo dos ingressantes do que no dos concluintes.
As notas máximas obtidas por concluintes ( 90,0 ) e ingressantes ( 84,5 ) foram diferentes.
As distribuições de notas para cada grupo se apresentaram com uma inclinação para notas maiores, como evidenciados pelos coeficientes de assimetria negativos.
A situação de assimetria é confirmada pelo Gráfico 3.5 apresentado, que proporciona uma comparação do desempenho entre concluintes e ingressantes em relação ao Componente Específico da prova.
Claramente a distribuição das notas médias dos ingressantes, à esquerda daquela dos concluintes, evidenciou o percentual mais elevado de notas médias nas faixas mais baixas para este grupo, com exceção do primeiro, que como já comentado, incluiu uma maior proporção de zeros entre os concluintes.
As medianas das distribuições das notas dos ingressantes e dos concluintes foram 49,4 e 54,8, respectivamente, enfatizando o melhor desempenho dos concluintes nesta parte da prova.
Assim como os Gráficos 3.2 e 3.4, o Gráfico 3.6 apresenta uma comparação dos resultados em relação às Regiões, Categoria Administrativa e Organização Acadêmica, agora levando em conta o desempenho dos ingressantes e dos concluintes no Componente Específico da prova.
Tendo como foco as notas médias dos estudantes, em determinada Região, observa -se que, entre os concluintes, as mais elevadas foram encontradas nas Regiões Sudeste ( 47,19 ) e Centro-Oeste ( 61,69 ), valores acima da média nacional para esta parte ( 46,27 ).
Entre os ingressantes, as maiores notas ocorreram também nas Regiões Nordeste ( 54,34 ), Sudeste ( 48,90 ) e Centro-Oeste ( 45,37 ), todas acima da média nacional ( 45,06 ).
A Região Sudeste apresenta os contingentes mais homogêneos com respeito à componente específica, tanto entre os ingressantes ( 1,18 ) como entre os concluintes ( 1,52 ).
Novamente, para a análise segundo a Categoria Administrativa, somente entidades federais, estaduais e particulares foram representadas.
Observa -se comportamento semelhante àquele da Formação Geral, ou seja, as médias mais elevadas, entre os concluintes foram encontradas em instituições de origem particular ( 55,34 ) e entre os ingressantes em instituições de origem federal ( 56,16 ).
Quanto à Organização Acadêmica, o comentário é o mesmo já feito para a nota bruta e referente à Formação Geral : somente as universidades foram representadas com este curso. 3.2 Análise das Questões Objetivas 3.2.1 Formação Geral A Tabela 3.4 apresenta as estatísticas básicas relativas às oito questões objetivas da Formação Geral dos estudantes.
Adiantamos que a primeira questão foi desconsiderada pelo critério da correlação ponto-bisserial.
Como pode ser observado na Tabela 3.4, os estudantes concluintes obtiveram desempenho inferior, com nota média de 50,32, em relação aos estudantes ingressantes cuja nota média foi 52,54.
A variabilidade entre os concluintes foi um pouco maior do que entre os ingressantes, como indicaram os desvios-padrão de 24,5, para o grupo de ingressantes e 29,4 para o grupo dos concluintes.
A amplitude total das notas foi a mesma para ambos os grupos : nota mínima de 0,00 e máxima de 85,7.
A mediana foi igual para ambos os grupos : 57,1.
A assimetria, medida pelo coeficiente de assimetria ( skewness ) foi um pouco maior em módulo para os ingressantes ( -0,7 ) do que para os concluintes ( -0,6 ), ainda que ambos fossem negativos indicando uma assimetria para a direita, fato reforçado pelos valores menores das médias dos dois grupos quando comparados à mediana correspondente.
A Tabela 3.5 apresenta a classificação das questões objetivas da Formação Geral, segundo o índice de facilidade ( questões em vermelho foram desconsideradas pelo critério ponto-bisserial ).
Como as questões foram comuns a todas as áreas, a tabela apresenta a classificação considerando -se os concluintes de todas as áreas e aqueles do curso de COMUNICAÇÃO SOCIAL – CINEMA.
De acordo com os índices obtidos, as questões objetivas da prova foram assim avaliadas para a área de COMUNICAÇÃO SOCIAL – CINEMA : das oito questões, três tiveram entre 60 \%, inclusive, e 85 \% de acertos e foram classificadas como fáceis ( somente uma para o total de concluintes ) ; outras duas foram consideradas medianas, alcançando entre 40 \%, inclusive, e 60 \% de acertos ( três para o total de concluintes ) ; duas obtiveram entre 15 \%, inclusive, e 40 \% de acertos, sendo consideradas difíceis ( três para o total de concluintes ).
Por fim, apenas uma foi classificada como muito difícil tento obtido percentual menor do que 15 \% de acertos ( mesmo número para o total de concluintes ), coincidentemente a questão desconsiderada pelo critério ponto-bisserial.
Nenhuma questão foi classificada na categoria muito fácil.
A Tabela 3.6, descreve os resultados das análises das questões objetivas relativas à Formação Geral, segundo o poder de discriminação, utilizando -se para tal o índice de discriminação ponto-bisserial.
Como as questões foram comuns a todas as áreas, a tabela apresenta a classificação considerando -se os concluintes de todas as áreas e aqueles do curso de COMUNICAÇÃO SOCIAL – CINEMA.
A desconsideração da questão segundo este critério considerou o universo de todos os concluintes.
Sete das oito questões apresentou índices iguais ou acima de 0,40 ( nenhuma para a população de concluintes COMUNICAÇÃO SOCIAL – CINEMA ), sendo, portanto, classificadas como muito boas nesta característica.
Para este grupo de estudantes, nenhuma foi classificada como boa com índice entre 0,30, inclusive, e 0,40 ( três para os concluintes como um todo ) ou classificada como média, com índice entre 0,20, inclusive, e 0,30 ( quatro para os concluintes COMUNICAÇÃO SOCIAL – CINEMA ).
Apenas uma das questões foi classificada como fraca para os concluintes de modo geral e também para os concluintes de COMUNICAÇÃO SOCIAL – CINEMA, sendo desconsiderada pelo critério. classificadas como fáceis ( somente uma para o total de concluintes ) ; outras duas foram consideradas medianas, alcançando entre 40 \%, inclusive, e 60 \% de acertos ( três para o total de concluintes ) ; duas obtiveram entre 15 \%, inclusive, e 40 \% de acertos, sendo consideradas difíceis ( três para o total de concluintes ).
Por fim, apenas uma foi classificada como muito difícil tento obtido percentual menor do que 15 \% de acertos ( mesmo número para o total de concluintes ), coincidentemente a questão desconsiderada pelo critério ponto-bisserial.
Nenhuma questão foi classificada na categoria muito fácil.
A Tabela 3.6, descreve os resultados das análises das questões objetivas relativas à Formação Geral, segundo o poder de discriminação, utilizando -se para tal o índice de discriminação ponto-bisserial.
Como as questões foram comuns a todas as áreas, a tabela apresenta a classificação considerando -se os concluintes de todas as áreas e aqueles do curso de COMUNICAÇÃO SOCIAL – CINEMA.
A desconsideração da questão segundo este critério considerou o universo de todos os concluintes.
Sete das oito questões apresentou índices iguais ou acima de 0,40 ( nenhuma para a população de concluintes COMUNICAÇÃO SOCIAL – CINEMA ), sendo, portanto, classificadas como muito boas nesta característica.
Para este grupo de estudantes, nenhuma foi classificada como boa com índice entre 0,30, inclusive, e 0,40 ( três para os concluintes como um todo ) ou classificada como média, com índice entre 0,20, inclusive, e 0,30 ( quatro para os concluintes COMUNICAÇÃO SOCIAL – CINEMA ).
Apenas uma das questões foi classificada como fraca para os concluintes de modo geral e também para os concluintes de COMUNICAÇÃO SOCIAL – CINEMA, sendo desconsiderada pelo critério.
O Gráfico 3.7 mostra, a título de ilustração, o comportamento das respostas à questão 7 de Formação Geral para os estudantes concluintes da área de COMUNICAÇÃO SOCIAL – CINEMA como função da nota geral na prova.
Trata -se de uma questão considerada de dificuldade fácil para os concluintes de COMUNICAÇÃO SOCIAL – CINEMA e de dificuldade média para o conjunto de concluintes de todas as áreas.
O poder de discriminação foi muito bom para os estudantes da área e bom para o conjunto de todos os concluintes.
Neste Gráfico, cada uma das cinco curvas, representou o percentual de respostas em determinada alternativa da questão, em função do escore total dos estudantes concluintes da área de COMUNICAÇÃO SOCIAL – CINEMA.
A curva azul marinho corresponde à alternativa A, a correta para esta questão.
Assim, observa -se que entre os estudantes com notas mais baixas – desempenho geral mais fraco na prova como um todo –, a tendência foi a escolha de alternativas incorretas, além da alternativa correta, especialmente a alternativa C – curva em vermelho seguida da alternativa E – curva em roxo.
A proporção de estudantes que selecionaram a resposta correta A, aumentou com o desempenho deste grupo de estudantes até os que obtiveram notas pelo menos 50,0, não chegando, porém, ao acerto de 100 \% dos estudantes.
Nos dois últimos intervalos apresenta um platô seguido de uma leve queda.
Esta análise permitiu verificar como a questão discriminou entre os grupos de desempenho.
Os Gráficos relativos às demais questões de Formação Geral constam do Anexo I. 3.2.2 Componente Específico A Tabela 3.7 apresenta as estatísticas básicas em relação às questões objetivas do Componente Específico.
Como pode ser observado, os estudantes concluintes obtiveram desempenho ligeiramente superior, com uma média de 48,56 pontos, em relação aos estudantes ingressantes com uma média de 47,80 pontos.
Os desvios-padrão indicaram que as variabilidade entre os concluintes foi maior ( 27,3 ) do que entre os ingressantes ( 21,1 ).
Os coeficientes de assimetria para os dois grupos foram negativos indicando uma inclinação para a direita.
A assimetria foi igual para ambos os grupos.
A análise das questões objetivas do Componente Específico em COMUNICAÇÃO SOCIAL – CINEMA, em relação ao grau de facilidade da prova ( ver Tabela 3.8 ), mostrou que, das 27 questões, oito encontraram -se na categoria difícil ( entre 15 \% e 40 \%, não incluso, de acertos ), 10 foram de nível médio ( entre 40 \% e 60 \%, não incluso, de acertos ), e 05 estavam na categoria fácil ( entre 60 \% e 85 \% de acertos ) e nenhuma teve a classificação muito fácil ( acima de 85 \% de acertos ).
Uma questão teve percentual de acertos menor do que 15 \% ( classificação muito difícil ) e três questões foram desconsideradas pela Comissão Assessora da Área de COMUNICAÇÃO SOCIAL – CINEMA.
Das 27 questões objetivas relativas ao Componente Específico ( ver Tabela 3.9 ), dezesseis obtiveram o índice de discriminação muito bom, quatro obtiveram o índice de discriminação bom e duas tiveram nível de discriminação médio.
O índice fraco, que eliminaria a questão, foi alcançado por duas questões e três questões foram canceladas pela comissão.
Na sequência, no Quadro 3.2 estão especificadas, por questão, as habilidades aferidas pelas questões do Componente Específico da prova, bem como os resultados da análise relativa aos índices de facilidade e de discriminação de cada questão.
Das 27 questões objetivas relativas ao Componente Específico ( ver Tabela 3.9 ), dezesseis obtiveram o índice de discriminação muito bom, quatro obtiveram o índice de discriminação bom e duas tiveram nível de discriminação médio.
O índice fraco, que eliminaria a questão, foi alcançado por duas questões e três questões foram canceladas pela comissão.
Na sequência, no Quadro 3.2 estão especificadas, por questão, as habilidades aferidas pelas questões do Componente Específico da prova, bem como os resultados da análise relativa aos índices de facilidade e de discriminação de cada questão.
O Gráfico 3.8 mostra, a título de ilustração, o comportamento da questão 11 do Componente Específico para os estudantes concluintes da área de COMUNICAÇÃO SOCIAL – CINEMA como função da nota geral na prova5.
Trata -se de uma questão considerada fácil para os concluintes de COMUNICAÇÃO SOCIAL – CINEMA.
O poder de discriminação foi muito bom para os estudantes da área.
Neste Gráfico, cada uma das cinco curvas representa o percentual de respostas em determinada alternativa da questão, em função da nota total dos estudantes concluintes de COMUNICAÇÃO SOCIAL – CINEMA.
A curva verde claro correspondeu à alternativa B, a correta para esta questão.
Assim, observou -se que entre os estudantes com escores mais baixos – desempenho geral mais fraco na prova como um todo –, a tendência foi a escolha de uma alternativa incorreta, a alternativa E – curva em roxo.
A proporção de estudantes que selecionaram a resposta correta B – curva verde claro – aumentou com o desempenho deste grupo de estudantes chegando ao acerto de 100 \% dos estudantes com notas pelo menos 75,0.
Esta análise permitiu verificar como a questão discriminou entre os grupos de desempenho.
Os Gráficos relativos às demais questões do Componente Específico constam do Anexo I. 3.3 Análise das Questões Discursivas 3.3.1 Formação Geral A Tabela 3.10 apresenta as estatísticas básicas da prova por grupo de estudantes.
A análise dos resultados de desempenho dos estudantes nas questões discursivas que são apresentados na tabela 3.10 e no gráfico 3.9, mostrou que as notas foram mais baixas no conjunto dessas questões que no das questões objetivas.
O total de estudantes inscritos para a prova foi de 967.
Destes, 74,7 \% compareceram, sendo que a presença foi maior entre concluintes ( 83,2 \% ) do que entre ingressantes ( 67,0 \% ).
A média geral da prova foi de 25,88, sendo que os ingressantes obtiveram média mais alta, 28,75, que os concluintes, cuja média foi 23,31.
O desvio-padrão geral foi de 35,2, sendo o do grupo de ingressantes maior ( 37,1 ), em relação ao grupo dos concluintes ( 33,2 ), aqueles possuíram uma distribuição menos homogênea nas notas.
A nota máxima foi 100,0, atingida por concluintes e ingressantes.
Os coeficientes de assimetria para os dois grupos considerados, o de concluintes e o de ingressantes, bem como para a população como um todo foram todos positivos, indicando uma maior concentração à esquerda da distribuição.
Para todos os grupos, mais de metade dos estudantes presentes obtiveram nota zero nas questões discursivas da Formação Geral.
Das 462 provas com nota zero, 26 ( 5,6 \% ) continham um protesto, 184 ( 39,8 \% ) estavam em branco e 252 ( 54,5 \% ) foram respondidas, porém de forma inadequada.
Entre os concluintes, a maioria ( 55,6 \% ) dentre aqueles com nota zero apresentou a prova em branco.
Entre os ingressantes, 53,3 \% das notas nulas decorreram de correção.
Uma comparação entre o desempenho de ingressantes e concluintes pode ser realizada a partir dos dados apresentados no Gráfico 3.9.
As notas dos concluintes, de um modo geral, foram mais baixas do que as dos ingressantes, cuja distribuição de notas está à direita dos primeiros.
A ideia de melhor desempenho foi reforçada pela proporção de notas nulas : 66,1 \% entre os concluintes e 61,6 \% entre os ingressantes.
Na sequência, discorrer -se-á, especificamente, sobre uma das questões discursivas de Formação Geral, a número 10 já que a número 9 foi desconsiderada.
Em especial, serão analisados os conteúdos englobados na questão e os resultados obtidos pelos estudantes.
Ao final, são apresentados os comentários da banca corretora a respeito do observado nas respostas dos estudantes, suas impressões e conclusões.
Cumpre esclarecer que, tendo em vista que as questões discursivas de Formação Geral foram padronizadas, ou seja, constam de todas as provas, os comentários da banca foram os mesmos para todas as áreas acadêmicas, eis que direcionados a todos os estudantes que participaram do ENADE/2009. 3.3.1.1 Análise da Questão Discursiva 9 de Formação Geral DESCONSIDERADA 3.3.1.2 Análise da Questão Discursiva 10 de Formação Geral A Questão Discursiva nº. 10 buscou avaliar as habilidades : • Ler e interpretar textos ; • Analisar e criticar informações ; • Extrair conclusões por indução e/ou dedução ; • Estabelecer relações, comparações e contrastes em diferentes situações ; • Fazer escolhas valorativas avaliando consequências ; • Questionar a realidade ; • Argumentar coerentemente.
O Quadro 3.3 apresenta as habilidades/competências, temas predominantes e estatísticas básicas da questão 10 de Formação Geral por grupo de estudantes.
Cumpre notar que como a questão 9 foi desconsiderada, a nota referente às questões discursivas da Formação Geral é idêntica à nota da questão 10.
A média geral da prova foi de 25,88, sendo que os ingressantes obtiveram média mais alta, 28,75, que os concluintes, cuja média foi 23,31.
O desvio-padrão geral foi de 35,2, sendo que o do grupo de ingressantes maior ( 37,1 ), em relação ao grupo dos concluintes ( 33,2 ), aqueles possuíram uma distribuição menos homogênea nas notas.
A nota máxima foi 100,0, atingida por concluintes e ingressantes.
Os coeficientes de assimetria para os dois grupos considerados, o de concluintes e o de ingressantes, bem como para a população como um todo foram positivos, indicando uma maior concentração à esquerda da distribuição.
Para todos os grupos, mais de metade dos estudantes presentes obtiveram nota zero nestas questões discursivas de Formação Geral.
Das 462 provas com nota zero, 26 ( 5,6 \% ) continham um protesto, 184 ( 39,8 \% ) estavam em branco e 252 ( 54,5 \% ) foram respondidas, porém de forma inadequada.
Entre os concluintes, a maioria ( 55,6 \% ) dentre aqueles com nota zero estava com a prova em branco.
Entre os ingressantes, 53,3 \% das notas nulas ( zero ) decorreram do desempenho do próprio estudante verificadas na correção.
Uma comparação entre o desempenho de ingressantes e concluintes para esta questão foi exatamente igual à já realizada para a nota discursiva de Formação Geral a partir dos dados apresentados no Gráfico 3.9. 3.3.1.3 Comentários sobre as respostas à Questão Discursiva 10 Apresentação A questão 10 de Formação Geral apresentou um questionamento capaz de analisar a habilidade do estudante em reconhecer o que é um direito social, optar por um direito social vigente que possa diminuir as desigualdades de renda familiar no Brasil e discutir os argumentos em favor do direito social escolhido como suporte à resposta.
Para análise das respostas apresentadas pelos estudantes, houve uma leitura inicial bastante criteriosa.
Foram considerados para efeito de pontuação, os seguintes argumentos : • Acesso à educação pública, gratuita e de qualidade, garantindo ao cidadão melhores condições de vida ; • Assistência à saúde, em seu contexto mais amplo, garantindo os cuidados essenciais de saúde, baseados em técnicas apropriadas ao alcance de todas as pessoas, contando com a participação da população ; • Permanência do estudante na escola, em todos os níveis escolares ( da educação infantil à educação superior, possibilitando a qualificação profissional ), o que garante melhores condições de trabalho e remuneração.
Considerou -se a qualificação e o desempenho docente, desenvolvendo nos estudantes melhores expectativas de aprendizagem e interesse pelo conteúdo da escola.
Os professores desenvolvendo metodologias adequadas e conhecimento pleno quanto a questões de aprendizagem ; • Condições dignas de trabalho, com remuneração que garanta qualidade de vida.
Foram consideradas as condições de trabalho, a preparação profissional para o trabalho e os cursos de aperfeiçoamento, oportunizando momentos de aprendizagem e desenvolvimento, a partir de situações desencadeadoras ; • Ser proprietário do imóvel em que reside, reduzindo os gastos com aluguel, garantindo melhores condições de vida, equilíbrio financeiro e construção do bem-estar familiar.
Desempenho dos estudantes na questão É importante salientar que havia respostas que desenvolveram adequadamente o questionamento apresentado, com fundamentação baseada em argumentos coerentes.
Contudo, em um número considerável de respostas, havia apenas um argumento e não dois, conforme fora solicitado.
Muitos estudantes informam sobre problemas da educação, de um modo geral.
Algumas \textbf{observações} mencionadas pelos corretores durante o processo de correção foram : • há preocupação com uma educação de qualidade, com sugestões de cursos profissionalizantes e métodos preparatórios de aulas. • a educação é revelada como de baixa qualidade, exaltando -se treinamentos e cursos de atualização para a classe docente. • a saúde como essência para a vida.
O SUS é colocado como parte integrante de direito social, garantindo às pessoas fácil acesso. • oportunidade de emprego como forma de ascensão social, direito, cidadania. • os políticos brasileiros como causadores do despreparo educacional.
As respostas analisadas apresentaram norma culta, coesão e coerência razoáveis, que garantiram a compreensão da mensagem desenvolvida pelos estudantes.
Os erros ortográficos foram poucos e não comprometeram a discussão sobre o assunto e sobre os pontos de vista abordados.
Quanto à coerência, houve um desvio quanto ao objetivo da questão, já que alguns estudantes, através dos seus argumentos retrataram reivindicações, reclamações e situações de problemas enfrentados pelos brasileiros, citando golpes políticos.
Houve erros quant a concordância verbal e nominal.
Houve, também, dificuldade em utilizar os conectores corretamente.
Alguns estudantes apresentaram em suas respostas redundância e prolixidade.
O desempenho foi considerado satisfatório, pois a maioria dos estudantes soube responder ao questionamento.
Houve domínio do conteúdo, pois as ideias, os argumentos e a originalidade corresponderam ao esperado. 3.3.2 Componente Específico Nas questões discursivas do Componente Específico – questões de números 38, 39 e 40 –, a diferença de desempenho entre os grupos não foi muito acentuada, inclusive menor do que para as questões discursivas de Formação Geral.
A Tabela 3.11 contém as estatísticas da média destas questões.
Observa -se que no grupo dos concluintes, a nota média alcançada foi 33,58, enquanto que no grupo dos ingressantes foi de 29,87, ambas maiores do que as obtidas nas questões discursivas de Formação Geral, que foram de 28,75 para os ingressantes e 23,31 para os concluintes ( Tabela 3.10 ).
As notas são menores dos que as obtidas para as questões objetivas desse componente ( Tabela 3.7 ) : respectivamente 47,80 e 48,56 para ingressantes e concluintes.
A análise de cada uma destas questões será feita nas subseções seguintes.
Os valores positivos para a medida de assimetria indicaram uma distribuição concentrada à esquerda, com menores valores.
Dos 722 presentes, 31,9 \% obtiveram nota zero.
A tabela 3.12 apresenta um detalhamento da situação destes estudantes. 50,7 \% não fizeram nenhuma das 3 questões.
Esta proporção foi menor entre os concluintes ( 40,8 \% ) do que entre os ingressantes ( 63,5 \% ).
A combinação de questões em branco e preenchidas correspondeu a 87,8 \% das provas.
Questões com protesto correspondeu a 2,7 \% do total das provas que receberam nota zero.
O Gráfico 3.10 apresenta a distribuição das notas nas questões discursivas do Componente Específico por grupo de estudantes.
A moda se situou no primeiro intervalo, aquele com valores entre 0 e 10 incluindo os extremos do intervalo.
A impressão de assimetria informada pelos coeficientes de assimetria positivos ( Tabela 3.11 ) foi reforçada no gráfico pela concentração no primeiro intervalo. 3.3.2.1 Análise da Questão Discursiva 38 do Componente Específico A questão 38, cujas habilidades aferidas encontram -se descritas no Quadro 3.4, que ainda contém outras informações, seguido da Tabela 3.13 e do Gráfico 3.11, foi aquela em que os grupos de estudantes observados obtiveram o melhor desempenho entre as questões discursivas do Componente Específico.
A nota média do grupo dos ingressantes foi de 45,36, e a dos concluintes 44,52.
O percentual de notas zero obtidas nos dois grupos foi o menor entre as questões do Componente Específico, de 34,2 \%, menor entre os ingressantes ( 30,8 \% ) do que entre os concluintes ( 37,3 \% ).
Quase 4/5 dos concluintes e dos ingressantes que tiraram zero deixaram a questão em branco ( ver Tabela 3.13 ).
Uma pequena minoria ( 9,3 \% ) utilizou a questão para protestar.
Apesar do baixo desempenho, em ambos os grupos a nota máxima de 100,0 pontos foi alcançada.
Os coeficientes de assimetria são próximos de zero indicando uma distribuição simétrica em torno da média.
O Gráfico 3.11 apresenta a distribuição das notas na questão 38 discursiva do Componente Específico por grupo de estudantes.
A moda se situou no primeiro intervalo, aquele com valores entre 0 e 10 incluindo os extremos do intervalo.
A impressão de assimetria informada pelos coeficientes de assimetria próximos de zero ( Quadro 3.4 ) foi enfraquecida no gráfico pela concentração no primeiro intervalo e concentrações crescentes nos intervalos subsequentes.
A distribuição apresentou outros máximos locais no último intervalo. 3.3.2.2 Comentários sobre as respostas à Questão Discursiva 38 A questão apresentou grau mediano de dificuldade, cobrando dos estudantes conhecimentos básicos sobre as mídias digitais e suas possibilidades e potencialidades.
Detectou -se, como resultado, um número alto de respostas superficiais, com a exposição de poucos conhecimentos sobre este novo espaço, o que considera -se contraditório à tendência apontada por este mercado e à realidade de rotina dos estudantes examinados ( afinal, o computador e o universo virtual são uma prática comum e necessária a todos os estudantes de nível superior do País ).
Exposições de conhecimentos sobre novas estéticas e experimentações que vêm sendo implantadas poderiam ser amplamente discutidas e apresentadas nesta questão, o que na grande maioria dos casos, não ocorreu. 3.3.2.3 Análise da Questão Discursiva 39 do Componente Específico A questão 39, cujas habilidades aferidas encontram -se descritas no Quadro 3.5, que ainda contém outras informações, seguido da Tabela 3.14 e do Gráfico 3.12, foi aquela em que os grupos de estudantes observados obtiveram desempenho intermediário entre as questões discursivas do Componente Específico.
A nota média do grupo dos ingressantes foi de 28,15, e a dos concluintes 30,33.
O percentual de notas zero obtidas nos dois grupos foi de 55,1 \% entre os ingressantes e 55,9 \% entre os concluintes.
Das provas com notas nulas nesta questão, quase 2/5 foram deixados em branco : 46,8 \% entre os ingressantes e 31,0 \% entre os concluintes ( ver Tabela 3.14 ).
Entre os com nota nula, uma pequena minoria ( 5,2 \% ) utilizou a questão para protestar.
Apesar do baixo desempenho, em ambos os grupos a nota máxima de 100,0 pontos foi alcançada.
Os coeficientes de assimetria positivos indicaram uma inclinação para á esquerda e cauda para a direita.
O Gráfico 3.12 apresenta a distribuição das notas na questão 39 discursiva do Componente Específico por grupo de estudantes.
A moda se situou no primeiro intervalo, aquele com valores entre 0 e 10 incluindo os extremos do intervalo.
A impressão de assimetria informada pelos coeficientes de assimetria positivos ( Quadro 3.5 ) foi reforçada no gráfico pela concentração no primeiro intervalo.
A distribuição não apresentou outros máximos locais. 3.3.2.4 Comentários sobre as respostas à Questão Discursiva 39 Embora a questão possuísse um grau mediano de dificuldade que exigiu dos estudantes conhecimentos técnicos e teóricos com relação ao cinema, a positiva participação destes pôde ser constatada nas respostas apresentadas.
Muitos estudantes discutiram com relativo grau de conhecimento sobre questões relativas à linguagem fílmica, aos gêneros, à história do cinema, às técnicas e aos meios de produção.
Pode -se considerar uma questão exclusivamente teórica que objetiva analisar o grau de conhecimento básico de um estudante acerca do cinema como técnica, e do filme como uma ferramenta de arte. 3.3.2.5 Análise da Questão Discursiva 40 do Componente Específico A questão 40, cujas habilidades aferidas encontram -se descritas no Quadro 3.6, que ainda contém outras informações, seguido da Tabela 3.15 e do Gráfico 3.13.
A questão foi aquela em que os grupos de estudantes observados obtiveram o pior desempenho entre as questões discursivas do Componente Específico.
A nota média do grupo dos ingressantes foi de 16,17, e a dos concluintes 25,94.
Nesta questão ocorre também o maior percentual de notas nulas : 73,3 \% entre os ingressantes e 61,9 \% entre os concluintes.
Aqui também a maioria das provas com nota zero decorreu de provas deixadas em branco : 67,6 \% entre os ingressantes e 53,0 \% entre os concluintes ( ver Tabela 3.15 ).
Uma pequena minoria ( 3,3 \% ), dentre aqueles com nota nula, utilizou a questão para protestar.
Em ambos os grupos, a nota máxima de 100,0 pontos foi alcançada.
Os coeficientes de assimetria positivos indicaram uma inclinação para a esquerda.
O Gráfico 3.13 apresenta a distribuição das notas na questão 40 discursiva do Componente Específico por grupo de estudantes.
A moda se situou no primeiro intervalo, aquele com valores entre 0 e 10 incluindo os extremos do intervalo.
A impressão de assimetria informada pelos coeficientes de assimetria positivos ( Quadro 3.6 ) foi reforçada no gráfico pela concentração no primeiro intervalo. 3.3.2.6 Comentários sobre as respostas à Questão Discursiva 40 A questão possui grau de dificuldade elevado ( principalmente por conter três questionamentos distintos, o que exigiu grande capacidade de síntese dos estudantes ao relacioná -los apenas nas dez linhas disponíveis ) e seu padrão de resposta não contemplou todos os questionamentos apresentados em seu enunciado.
O padrão de resposta apresenta abordagens às propostas e características do Cinema Novo e sua projeção no cinema atual, mas não faz nenhuma menção sobre suas propostas estéticas com o pensamento político, omitindo o fato de que esteticamente ocorreu uma ruptura com o cinema considerado “ comercial ”, demonstrando que o movimento vinha em oposição a todo o padrão estabelecido, inclusive o político em que os equívocos da política desenvolvimentista e da ditadura eram constantemente criticados.
Poucos estudantes contemplaram em suas dissertações os três itens propostos, focando o último com mais intensidade por se encontrar mais próximo de suas realidades, conseguindo expor, ainda que de forma superficial, correlações importantes entre o Cinema Novo e a atual produção cinematográfica brasileira.
As questões vinculadas à habilitação Cinema foram as que apresentaram maior qualificação nas respostas entregues para correção.
Muitos estudantes mostraram -se acadêmica e socialmente mais preparados para o exercício do pensamento crítico e original, expressos tanto na abordagem que analisa produção de sentido, quanto nas especificidades técnicas do campo.
Há poucas instituições voltadas para a atividade, e o afunilamento à especialização manifestam -se através da maior sofisticação dos textos dos estudantes e do conhecimento do contexto.
Aparentemente, todos os campos básicos foram contemplados nas perguntas discursivas de Cinema : a passagem de técnica analógica para a digital, a identificação dos elementos que compõem uma análise fílmica e o conhecimento da história do cinema nacional como referência para a identificação das características da produção contemporânea.
Erros de gramática e de grafia foram poucos especialmente entre as respostas com estrutura adequada de texto.

% matched lemmas: observação
\end{textsample}
